\subsubsection{4.2 Die Lücke zum gesamten übrigen
Raum}\label{die-luxfccke-zum-gesamten-uxfcbrigen-raum}

Die native, nicht nur perturbative Casimiridentität lautet

\[\adjustbox{max width=\linewidth}{$\displaystyle 
 \sum_AP_A^\dagger P_A=\tfrac12(15N_f-C_{\rm Spin(10)}-C_{\rm SU(4)}).
$}\]

In einem Block \(N=n,N_b=b\) ist die rechte Seite höchstens
\(\tfrac{15}2(n-2b-\eta)\) mit \(\eta=n\bmod2\). Bei ungerader
Fermionzahl kommt der Casimir-Mindestwert 15 hinzu. Die Norm des
Boson-Erzeugungszeilenoperators auf dem Zielblock ist \(\sqrt{b+1}\).
Damit liefert die Blocknorm-Abschätzung eine skalare Jacobi-Untergrenze
mit Diagonale \(b\) und quadrierten Nebendiagonalen

\[\adjustbox{max width=\linewidth}{$\displaystyle 
 a_b^2=\frac{15(b+1)(n-2b-\eta)}{800}
$}\]

in Einheiten \(\Delta=1\). Man nimmt negative Nebendiagonalen; die
Quadratformabschätzung benutzt die Normen der vollen
Bosonzahlkomponenten eines beliebigen Zustands, nicht einzelne
Testvektoren.

Für \textbf{jedes \(n=65,\ldots,127\)} wurden sämtliche LDL-Pivots der
Vergleichsmatrix oberhalb \(-4/5\) exakt positiv nachgewiesen. Der erste
Bosonindex ist \(\lceil(n-64)/2\rceil\), der letzte
\(\lfloor n/2\rfloor\). Somit gilt \(H_n>-.8\Delta\) in allen 63 hohen
Ladungssektoren.

Für alle \(n\le62\) liefert die Minimierung der Cauchy-Schwarz-Schranke

\[\adjustbox{max width=\linewidth}{$\displaystyle 
 H_n/\Delta\ge-\frac{n-\eta}{4}(\sqrt{23/20}-1)>-1.121899.
$}\]

Jede andere Ladungsverteilung als (63,64) oder (64,63) enthält eine Bank
mit \(n\le62\) und eine mit \(n\ge65\). Ihre Gesamtenergie liegt also
oberhalb \(-1.921899\Delta\). Die elektrische Energie ist nichtnegativ.
Das obere Bandende von \(P\) ist kleiner als
\(-2.225448\Delta+\kappa/2\).

Für die beiden zentralen Ladungsverteilungen wurden außerdem die
Vergleichsmatrizen ab \(b=1\) bei \(n=63,64\) oberhalb \(-3/4\) erneut
rational geprüft. Durch Minmax mit Kodimension 64 beziehungsweise eins
ergeben sich die bekannten angeregten Energieböden; es wird nicht
behauptet, dass der exakte Spektralprojektor mit dem
Null-Bosonenprojektor identisch ist.

Zusammen ergibt sich eine Lücke zwischen \(P\) und \(Q=I-P\) von
mindestens

\[\adjustbox{max width=\linewidth}{$\displaystyle 
 \boxed{\delta=.303549\Delta-\kappa/2>0.}
$}\]

Die anderen beiden Vergleichslücken sind \(.345812\Delta-\kappa/2\) und
\(.379636\Delta-\kappa/2\), also größer. Diese Schranke umfasst
\textbf{alle} Ladungsaufteilungen, Bosonzustände und internen Moden des
deklarierten physikalischen Sektors.

\subsubsection{4.3 Gleichmäßige Auslaufkontrolle und endliche
Transferzeit}\label{gleichmuxe4uxdfige-auslaufkontrolle-und-endliche-transferzeit}

Setze \(v=64|\tau|\) und fordere \(2v<\delta\). Sei \(a\) das obere
Eigenwertende von \(H_0|_P\). Minmax liefert ein Spektralband \(P'\) von
\(H\) derselben Dimension 128, dessen oberer Rand ≤\(a+v\) ist. Die
Kompression \(QHQ\) liegt ≥\(a+\delta-v\).

Mit \(Y=QP'\), als Abbildung aus \(\operatorname{Ran}P'\), gilt die
Sylvestergleichung

\[\adjustbox{max width=\linewidth}{$\displaystyle 
 (QHQ)Y-Y(H|_{P'})=-QVP\,PP'.
$}\]

Die beiden Spektren sind geordnet und mindestens \(\delta-2v\) getrennt.
Die Lösung über das konvergente Exponentialintegral liefert
\(\|Y\|\le v/(\delta-2v)\): nach einem gemeinsamen Skalarshift ist der
Integrand durch \(v e^{-s(\delta-2v)}\) beschränkt. Wegen gleicher
endlicher Ränge gilt dieselbe Schranke für \(\|P'-P\|\).

Da \(P'\) mit \(H\) kommutiert, folgt zu \textbf{jeder} Zeit

\[\adjustbox{max width=\linewidth}{$\displaystyle 
 \boxed{\|Qe^{-itH}P\|\le L:=\frac{2v}{\delta-2v}.}
$}\]

Aus der projizierten Schrödingergleichung und Duhamel folgt weiter

\[\adjustbox{max width=\linewidth}{$\displaystyle 
 \|Pe^{-itH}\psi-e^{-itPHP}\psi\|
 \le |t|vL=:\eta(t),\qquad \psi\in P,\ \|\psi\|=1.
$}\]

Für den Zielprojektor \(P_R\subset P\) weichen die Wahrscheinlichkeiten
damit höchstens um \(2\eta(t)\) ab. Dies ist eine nichtperturbative
Fehlergrenze für die volle Entwicklung; nur der ausgewählte
Kopplungsbereich ist klein. Es handelt sich nicht um das Weglassen eines
unbekannten höheren Störungsterms.

\subsubsection{4.4 Ein vollständig numerisch spezifizierter, rational
zertifizierter
Punkt}\label{ein-vollstuxe4ndig-numerisch-spezifizierter-rational-zertifizierter-punkt}

Mit \(\hbar=1\) wähle

\[\adjustbox{max width=\linewidth}{$\displaystyle 
 \tau/\Delta=10^{-8},\quad \kappa/\Delta=10^{-10},
 \quad T:=t\Delta=169500000.
$}\]

Diese Zahlen sind ein konservativer Existenzpunkt, \textbf{keine aus
TFPT abgeleiteten Konstanten und keine Behauptung schneller oder
optimaler Übertragung}. Die physikalische Sekundenskala ist nicht
bestimmt.

Die genaue komprimierte Rabi-Wahrscheinlichkeit ist

\[\adjustbox{max width=\linewidth}{$\displaystyle 
 p_P(t)=\frac{|\tau|^2Z^2}{|\tau|^2Z^2+(\kappa/4)^2}
 \sin^2\!\left(t\sqrt{|\tau|^2Z^2+(\kappa/4)^2}\right).
$}\]

Für \textbf{jedes} zulässige unbekannte \(Z\) liegt der Winkel weniger
als 0,08 von \(\pi/2\) entfernt. Das folgt aus
\(Z_{\rm lo},Z_{\rm hi}\), \(\sqrt{x^2+y^2}\le x+y^2/(2x)\) und
rationalen Pi-Grenzen. Letztere wurden zusätzlich mit der
Machin-Identität und endlichen alternierenden Arctan-Reihen
eingeschlossen; Fließkomma-Pi ist keine Beweisvoraussetzung.

Mit \(\sin^2(\pi/2+u)\ge1-u^2\) folgt

\[\adjustbox{max width=\linewidth}{$\displaystyle 
 p_P(T)>\left(1-\frac{\kappa^2}{16\tau^2Z_{\rm lo}^2}\right)
 (1-.08^2)>.993591982278.
$}\]

Die vollständigen Fehlergrenzen sind

\[\adjustbox{max width=\linewidth}{$\displaystyle 
 L=\frac{25600}{6070954399}<4.216800\cdot10^{-6},
 \quad L^2<1.8\cdot10^{-11},
 \quad\eta(T)=\frac{2777088}{6070954399}<.000457438455.
$}\]

Also

\[\adjustbox{max width=\linewidth}{$\displaystyle 
 \boxed{p_{\rm voll}(T)>p_P^{\rm lo}-2\eta(T)
 =\frac{25218291754002904578202303151726961}
 {25404324948782158166703488466197500}>.992677105368.}
$}\]

Die unbekannten \(E_0,E_h\) treten nur als gemeinsame Phase auf und
müssen für diese Schranke nicht genau ausgerechnet werden. Eine arbiträr
hohe Treue ist innerhalb des zusätzlichen Modells prinzipiell
erreichbar: bei bekanntem \(Z\), \(\kappa/|\tau|\to0\),
\(|\tau|/\Delta\to0\) und \(t\sim\pi/(2|\tau|Z)\) verschwinden sowohl
Detuning als auch der Fehler \(tv^2/\delta\). Das ist eine bedingte
Grenzaussage mit wachsender Laufzeit, kein ausführbares natives
Optimierungsverfahren.

\subsubsection{4.5 Dieselbe ursprüngliche Fermionantwort ohne
anfänglichen
Filter}\label{dieselbe-urspruxfcngliche-fermionantwort-ohne-anfuxe4nglichen-filter}

Für \(\psi_f=(f_{x,r}\Omega_x/\sqrt\nu)\otimes\Omega_y\otimes|0\rangle\)
gilt

\[\adjustbox{max width=\linewidth}{$\displaystyle 
 \psi_f=\sqrt w\,|L,r\rangle+\sqrt{1-w}\,q,
 \quad q\in Q,\quad w=Z/\nu\ge Z_{\rm lo}/Z_{\rm hi}.
$}\]

Die Auslaufnorm in umgekehrter Richtung erfüllt ebenfalls
\(\|Pe^{-itH}Q\|\le L\), indem man die obige Schranke adjungiert und die
Zeit umkehrt. Der unerwünschte Anteil kann also nicht beliebig stark
destruktiv in das niedrige Zielband einstreuen. Mit der
Dreiecksungleichung, \(2\sqrt{w(1-w)}\le1\) und der gesondert geprüften
Positivität vor dem Quadrieren folgt

\[\adjustbox{max width=\linewidth}{$\displaystyle 
 \boxed{\|P_Re^{-itH}\psi_f\|^2
 \ge w\,p_{\rm voll}^{\rm lo}-L
 >.897260353455>.897.}
$}\]

Damit ist \textbf{kein vorbereitender Projektor auf die isolierte Linie}
nötig, um eine starke Übertragungsaussage über die ursprüngliche
Entnahmeantwort zu erhalten. Noch benötigt werden die deklarierte
geladene Präparation und der zusätzliche Link. Die Aussage betrifft das
niedrige Zielband; sie behauptet nicht, dass die vollständige
hochenergetische Restantwort ebenfalls formtreu übertragen wird. Für
diesen ungefilterten Anfangszustand gilt die extrem kleine
Auslaufwahrscheinlichkeit aus 4.3 nicht: Er startet bereits teilweise
außerhalb von \(P\).

\subsection{5. Warum die Herkunftslücke nicht durch längeres Rechnen
verschwindet}\label{warum-die-herkunftsluxfccke-nicht-durch-luxe4ngeres-rechnen-verschwindet}

Sei \(\Pi_x=(-1)^{N_{f,x}}\) die lokale Fermionparität. Die bisher
angegebenen bankinternen Paarumwandlungen, Zahloperationen,
zahlenerhaltenden Symmetrielifts sowie reine Bosonverbindungen
kommutieren mit jeder \(\Pi_x\). Jede endliche Komposition, lineare
Kombination, Adjungierung und jeder durch solche Operatoren realisierte
einzelne Messzweig tut das ebenfalls. Starke beschränkte Grenzwerte
bleiben im Kommutanten der Parität.

Der benötigte Link erfüllt dagegen \(\Pi_xB=-B\Pi_x\) und
\(\Pi_yB=-B\Pi_y\), während er die Gesamtparität erhält. \textbf{Er
liegt nicht in dieser bekannten lokalen geraden Algebra.} Auch adaptives
Wiederholen und Nachselektieren gerader Krauszweige kann die
Sektorgrenze nicht überwinden.

Die Voraussetzung „jeder realisierte Zweig ist gerade`` ist wesentlich.
Eine bloß paritätskovariante CP-Abbildung kann ungerade Krausoperatoren
besitzen. Der im Prüfer enthaltene Reset-Kanal auf einer Fermionmode ist
ein exaktes Gegenbeispiel zur unzulässigen stärkeren Behauptung.

Das beweist keinen universellen Unmöglichkeitssatz für den Compiler. Es
beweist, \textbf{welche neue Quelleneigenschaft gesucht werden muss}:
eine gerade Gesamtoperation, die bezüglich der zwei operational
definierten Teilbereiche ungerade ist, oder eine ursprüngliche
Überlappung, durch die die angenommene Zerlegung in unabhängig gerade
Banken falsch war. Eine andere bloße Clockphase oder weitere
Boson-Paarumwandlung reicht nicht.

\subsection{6. Was von der arithmetischen Untersuchung
trägt}\label{was-von-der-arithmetischen-untersuchung-truxe4gt}

\subsubsection{6.1 Eine einzelne endliche Bank ist nicht der volle
Logarithmusgenerator}\label{eine-einzelne-endliche-bank-ist-nicht-der-volle-logarithmusgenerator}

Quadratisches Ergänzen liefert mit der nativen Casimirnorm
\(\|\sum P_A^\dagger P_A\|\le480\)

\[\adjustbox{max width=\linewidth}{$\displaystyle 
 H_{\rm nat}\ge\tfrac\Delta2N_b-C,
 \qquad C=960g^2/\Delta.
$}\]

Die externe Rechnung verwendet teilweise die schwächere Dreiecksgrenze
\(C=7680g^2/\Delta\), die ebenfalls gültig ist. Über Minmax,
\textbf{nicht} eine allgemeine Operator-Monotonie der
Exponentialfunktion, folgen

\[\adjustbox{max width=\linewidth}{$\displaystyle 
 N_H(E)\le2^{64}\binom{\lfloor2(E+C)/\Delta\rfloor+60}{60},
 \qquad \operatorname{Tr}e^{-\beta H}
 \le\frac{2^{64}e^{\beta C}}{(1-e^{-\beta\Delta/2})^{60}}.
$}\]

Das Zustandswachstum ist polynomial, während ein vollständiges Spektrum
\(E_*\log n+E_{\rm off}\), \(E_*>0\), exponentiell viele Zustände unter
wachsender Energie verlangt. Die ursprüngliche Bank kann dieses
komplette Spektrum daher nicht energieerhaltend mit nur affiner
Skalenänderung tragen. Diese Schranke widerlegt nicht jeden
arithmetischen Untersektor, jede nichtlineare Umparametrisierung oder
jeden Nullstellenoperator.

Endlich viele solche Banken und endlich viele energetisch kontrollierte
Rotoren ändern den qualitativen Gegensatz nicht. Unendlich viele
identische Banken lösen ihn nicht automatisch: Bei gleichmäßig
beschränkten lokalen Anregungskosten divergiert bereits die globale
Wärmespur. Ein geeigneter relativer oder lokaler Spurbegriff wäre ein
zusätzlicher Gegenstand.

\subsubsection{6.2 Primzahlen können eine gewählte Dynamik organisieren,
erzwingen sie aber
nicht}\label{primzahlen-kuxf6nnen-eine-gewuxe4hlte-dynamik-organisieren-erzwingen-sie-aber-nicht}

Auf \(\ell^2(\mathbb N)\) kann man
\(H_{\rm ar}|n\rangle=\log n|n\rangle\) definieren. Die additive
Primfaktorbesetzung erklärt dann die Eulerstruktur. Der
Definitionsbereich lautet \(\sum_n(\log n)^2|\psi_n|^2<\infty\). Dass
quantenstatistische Systeme mit Zeta-Zustandssumme existieren, ist
bereits ein Ergebnis von
\href{https://alainconnes.org/wp-content/uploads/bostconnesscan.pdf}{Bost
und Connes}. Das ist ein relevanter Anschluss, nicht die fehlende native
Herleitung.

Der gelieferte Cocycle \(c(m,n)=(-1)^{v_2(m)v_3(n)}\) ist assoziativ,
weil Bewertungen unter Multiplikation additiv sind. Für
\(T_m|n\rangle=c(m,n)|mn\rangle\) gilt \(T_lT_m=c(l,m)T_{lm}\),
insbesondere \(T_2T_3=-T_3T_2\). Die unkontrollierten CP-Kanäle
verlieren diese globale Phase und erfüllen
\(\mathcal E_l\mathcal E_m=\mathcal E_{lm}\). In kohärent kontrollierten
Wegen bleibt die relative Phase messbar. Das ist eine exakte
Phasenstruktur, aber kein Beweis einer Quellenauswahl dieses Cocycle
oder von RH.

Strikt multiplikative Schritte \(m>1\) ergeben auch keine geschlossenen
Bahnen: \(n\mapsto mn>n\). Die Frequenz \(\log p\) eines Operators, eine
Bahnlänge \(r\log p\) und der Schwingungszeitraum \(2\pi/\log p\) müssen
auseinandergehalten werden. Hier hat der RH-Graph-Skill die Prüfung auf
Generator, Spur und Phase getrennt und eine bloße Analogieschließung
verhindert.

\subsubsection{6.3 Der RH-Zielsatz bleibt
unbewiesen}\label{der-rh-zielsatz-bleibt-unbewiesen}

Die vorgeschlagene Identität

\[\adjustbox{max width=\linewidth}{$\displaystyle 
 \frac{\Xi(z)}{\Xi(0)}=\det(I-z^2A^{-2}),\quad
 \Xi(z)=\xi(\tfrac12+iz),
$}\]

für einen strikt positiven selbstadjungierten Operator \(A\) mit
kompakter Inverser und \(A^{-2}\) von Spurklasse wäre tatsächlich
hinreichend: Die gesamte rechte Funktion hat ihre Nullstellen nur bei
reellen \(\pm\lambda_j(A)\), einschließlich Multiplizitäten. Aber dieser
Operator und die gesamte Funktionsidentität sind \textbf{nicht
konstruiert}. Ein Operator mit lediglich den Imaginärteilen der
Nullstellen wäre unzureichend.

Weder diese Prüfung noch die im externen Code bereits verwendete
Faktorisierung kleiner Testzahlen liefert einen neuen effizienten
Faktorisierungsalgorithmus, eine native Quantenoperation dafür oder eine
Lösung von P versus NP. Hylæan wurde in dieser Revision nicht neu
geprüft.

\subsection{7. Die lokale Polnäherung wird präziser, das relativistische
Feld bleibt
offen}\label{die-lokale-polnuxe4herung-wird-pruxe4ziser-das-relativistische-feld-bleibt-offen}

Die externe Resolventenabschätzung ist korrekt: Im Kreis mit Radius
\(.1\Delta\) um den \textbf{tatsächlichen}, unbekannten negativen Pol
\(-\epsilon\) gilt

\[\adjustbox{max width=\linewidth}{$\displaystyle 
 G(z)=\frac{Z}{z+\epsilon}+R(z),\qquad
 |R(z)|<.505213/\Delta.
$}\]

Der übrige Spektralträger ist mindestens \(.337373\Delta\) vom Pol
entfernt. Die relative Abweichung vom echten Polterm beträgt im Kreis
unter 5,75 \%. Die Abschätzung setzt keine Intervallmitte anstelle von
\(Z,\epsilon\) ein und ist keine globale Vernachlässigbarkeit der
Selbstenergie.

Aus v1.6.4 bleibt bestehen: Die exakte Zwei-Feld-Resolvente besitzt
einen nichtverschwindenden Rest \(\Sigma_2(z)\). Seine komplette
Funktion ist noch nicht bestimmt. Das neue kontrollierte Transportlemma
benötigt sie nicht vollständig, weil es den Rest durch eine
Spektrallücke kontrolliert.

Ebenso bleibt der Feldtypentest bestehen: Der native antisymmetrische
Paartensor zusammen mit zwei gleichhändigen Weylfeldern und einem
skalaren Vermittler ergibt den verschwindenden Kanal. Ein symmetrischer
Lorentz- Spinortensor oder eine zusätzliche gleichgeladene
zweidimensionale Antisymmetriemarke eröffnet einen nichtverschwindenden
Kanal, führt aber zusätzliche Struktur ein. Ein internes
Spin(10)-Spinorlabel ist kein Lorentz-Weylindex. Der neue Link behebt
diese Typfrage nicht. Die Konventionen des bisherigen Tests sind mit der
Primärübersicht von \href{https://arxiv.org/abs/0812.1594}{Dreiner,
Haber und Martin} verbunden; ein neuer vollständiger relativistischer
Adapter wurde hier nicht konstruiert.

Die spätere Feldtypenprüfung in \codeword{LATE\_AUDIT.md} zeigt außerdem
einen Lorentz-erlaubten Zwei-Ableitungs-Kandidaten für den symmetrischen
Spinortensor und widerlegt einen pauschalen (1,1)-Ausschluss: Der
Energie-Impuls-Tensor ist schon bei Dimension vier zulässig. Weder ein
gesunder nativer Feldadapter noch ein dynamischer Gravitonpol folgt
daraus.

